% LAMP-tau User Manual -- LaTeX edition.
% Build with: xelatex manual.tex (twice, for the table of contents)
% The English edition is generated from docs/manual.md; the Markdown
% edition also contains a Chinese summary.
\documentclass[11pt,a4paper]{article}
\usepackage[margin=2.2cm]{geometry}
\usepackage{fontspec}
\usepackage{amsmath}
\usepackage{amssymb}
\usepackage{booktabs}
\usepackage{longtable}
\usepackage{array}
\usepackage{xcolor}
\usepackage{listings}
\usepackage{textcomp}

\lstset{
  basicstyle=\ttfamily\small,
  breaklines=true,
  columns=fullflexible,
  keepspaces=true,
  frame=single,
  backgroundcolor=\color{gray!8},
  showstringspaces=false,
  extendedchars=true
}

\providecommand{\tightlist}{%
  \setlength{\itemsep}{0pt}\setlength{\parskip}{0pt}}
\providecommand{\texorpdfstring}[2]{#1}
\providecommand{\linethickness}{0.4pt}

\title{LAMP-tau User Manual}
\author{LAMP-tau contributors}
\date{Version 1.0.0}

\begin{document}
\maketitle
\tableofcontents
\newpage

\begin{quote}
\textbf{Abstract.} This manual covers installation, a runnable DyAC example, the full spin-phonon-coupling workflow, the YAML and file formats, output files, the Python API, HPC notes and troubleshooting. The Markdown edition in \texttt{docs/manual.md} retains a Chinese summary.
\end{quote}

Version 1.0.0 --- accompanying the manuscript \emph{Accelerating ab
initio spin-phonon relaxation simulation of single-ion magnets by
quantum embedding and spatial truncation} (J. Chem. Phys., 2026).

\begin{center}\rule{0.5\linewidth}{0.4pt}\end{center}

\section{Overview}\label{overview}

LAMP-tau is a Python toolkit for \emph{ab initio} spin-phonon relaxation
simulation of single-ion magnets (SIMs). It has two parts:

\begin{itemize}
\tightlist
\item
  \textbf{\texttt{mag\_anal}} --- analysis of electronic-structure
  results: crystal-field parameters (CFPs), effective spin Hamiltonian,
  \emph{g} tensors, spin-phonon coupling (SPC) matrix elements and
  relaxation times \(\tau\)(T).
\item
  \textbf{\texttt{SPC\_Tool}} --- scripts that prepare and organise the
  underlying calculations: VASP phonons, displaced structures,
  DMET+CASSCF-SO (PySCF) or CASSCF (OpenMolcas) inputs, and CFP
  extraction.
\end{itemize}

The electronic-structure workflow is

\begin{lstlisting}
VASP phonons ----> displaced structures ----> DMET+CASSCF-SO / CASSCF ----> CFPs
                                                                      |
                                                     SPC parameters <---+
                                                            |
                                                            ▼
                                                         tau(T), Ueff, tau0
\end{lstlisting}

\texttt{mag\_anal} is the part used in step 4--5; \texttt{SPC\_Tool}
covers steps 1--3.

\subsection{Requirements at a
glance}\label{requirements-at-a-glance}

\begin{longtable}[]{@{}>{\raggedright\arraybackslash}p{0.30\linewidth} >{\raggedright\arraybackslash}p{0.66\linewidth}@{}}
\toprule
Item\strut
 & Requirement\strut
\tabularnewline
\midrule
\endhead
Python\strut
 & \(\geq\) 3.10 (the code uses \texttt{match}/\texttt{case})\strut
\tabularnewline
Core Python packages\strut
 & NumPy, SciPy, PyYAML, prettytable, Matplotlib, ASE, SymPy\strut
\tabularnewline
DMET+CASSCF-SO\strut
 & PySCF, \texttt{basis\_set\_exchange},
\texttt{LAMP-emb} (\texttt{https://github.com/ZGuan97/LAMP\_emb})
(\texttt{embed\_sim})\strut
\tabularnewline
Phonons\strut
 & VASP (periodic) or a molecular Hessian (e.g.~ORCA)\strut
\tabularnewline
Alternative electronic structure\strut
 & OpenMolcas (for the Molcas path)\strut
\tabularnewline
Scheduler\strut
 & Slurm, PBS or Tianhe; the shipped \texttt{.sh} files are templates\strut
\tabularnewline
\bottomrule
\end{longtable}

VASP, OpenMolcas and ORCA are external programs and are \textbf{not}
distributed with LAMP-tau.

\begin{center}\rule{0.5\linewidth}{0.4pt}\end{center}

\section{Installation}\label{installation}

\subsection{From a clone}\label{from-a-clone}

\begin{lstlisting}
git clone https://github.com/ccme-tmc/LAMP_tau.git
cd LAMP_tau
python -m pip install -e .
\end{lstlisting}

This installs the \texttt{mag\_anal} package and the \texttt{lamp-tau}
command. If you only need the analysis part:

\begin{lstlisting}
python -m pip install -r requirements.txt
\end{lstlisting}

\subsection{LAMP-emb and PySCF}\label{lamp-emb-and-pyscf}

The DMET part uses the \texttt{embed\_sim} package from LAMP-emb:

\begin{lstlisting}
git clone https://github.com/ZGuan97/LAMP_emb.git
export PYTHONPATH=/path/to/LAMP_emb:$PYTHONPATH
python -m pip install pyscf basis_set_exchange
\end{lstlisting}

\texttt{dmet\_tmp.py.tmpl} imports \texttt{embed\_sim.myavas},
\texttt{ssdmet}, \texttt{siso}, \texttt{sacasscf\_mixer} and
\texttt{rdiis}.

\subsection{Check the installation}\label{check-the-installation}

\begin{lstlisting}
lamp-tau --help
python -c "import mag_anal; print(mag_anal.__version__)"
\end{lstlisting}

On a headless machine set:

\begin{lstlisting}
export MPLBACKEND=Agg
\end{lstlisting}

\subsection{Optional development
tools}\label{optional-development-tools}

\begin{lstlisting}
python -m pip install -e ".[dev]"   # pytest, ruff
python -m pytest tests -q
\end{lstlisting}

\begin{center}\rule{0.5\linewidth}{0.4pt}\end{center}

\section{Repository layout}\label{repository-layout}

\begin{lstlisting}
LAMP_tau/
|-- mag_anal/
|   |-- Anal_Tau.py          # tau(T) analysis; CLI entry point
|   |-- psh.py               # pseudospin Hamiltonian, g tensors, CFP projection
|   |-- gen_structures.py    # displaced structures from phonon modes
|   |-- Select_group.py      # molecular-group selection / spatial truncation
|   |-- dictionary.py        # element/mass dictionaries
|-- SPC_Tool/
|   |-- inpfile/             # input templates (VASP, PySCF, OpenMolcas)
|   |-- opt_freq/vaspsh/     # VASP geometry optimisation and frequencies
|   |-- Per_B_kp/Pyscfsh/    # displaced structures, inputs, CFP extraction
|   |-- sh/                  # scheduler templates
|   |-- Relaxation/          # legacy example configuration
|   |-- manual/              # Chinese manual (LaTeX + PDF)
|-- examples/relaxation/     # runnable DyAC example
|-- docs/                    # workflow, formats, HPC and data notes
|-- tests/                   # smoke test
|-- pyproject.toml
|-- requirements.txt
|-- CITATION.cff
|-- LICENSE
\end{lstlisting}

\begin{center}\rule{0.5\linewidth}{0.4pt}\end{center}

\section{Quick start: relaxation analysis
(DyAC)}\label{quick-start-relaxation-analysis-dyac}

The shipped example uses DyAC and runs from \texttt{examples/relaxation}.

\begin{lstlisting}
cd examples/relaxation
lamp-tau --yml inpfile/CFPs_v1.yml
\end{lstlisting}

or, without installing:

\begin{lstlisting}
PYTHONPATH=../.. python -m mag_anal.Anal_Tau --yml inpfile/CFPs_v1.yml
\end{lstlisting}

What happens:

\begin{enumerate}
\def\labelenumi{\arabic{enumi}.}
\tightlist
\item
  \texttt{Read\_input} parses the YAML file and creates
  \texttt{Out\_foldname}.
\item
  \texttt{Read\_freq} reads \texttt{freq.log}/\texttt{POSCAR} and builds
  the phonon modes.
\item
  \texttt{Read\_CFPs} reads \texttt{CFPs\_py.out} and fits ∂Bkq/∂Q.
\item
  \texttt{Read\_Static\_effH} reads the equilibrium output
  (\texttt{eqpyout}) and the tensor-operator matrices.
\item
  \texttt{Analyse\_Relax} computes the spin-phonon matrix, transition
  probabilities and \(\tau\)(T).
\end{enumerate}

Expected console tail:

\begin{lstlisting}
TestU_{eff} of Orbach is 190.1 K
log(tau) is -11.30 s
rate-determing step
\end{lstlisting}

Outputs (in \texttt{examples/relaxation/TPs\_file/} by default):

\begin{longtable}[]{@{}>{\raggedright\arraybackslash}p{0.30\linewidth} >{\raggedright\arraybackslash}p{0.66\linewidth}@{}}
\toprule
File & Content\tabularnewline
\midrule
\endhead
\texttt{OB\_v1.dat} & \texttt{1/T}, \texttt{-log10(\(\tau\)/s)}, standard
deviation\tabularnewline
\texttt{TPs1.png} & \texttt{log10(\(\tau\)/s)} vs \texttt{1/T}
(Orbach)\tabularnewline
\texttt{TPs2.png} & comparison figure\tabularnewline
\texttt{VP.log} & spin-phonon matrix elements\tabularnewline
\texttt{TPs/*.tps}, \texttt{TPs/*.tp} & cached transition matrices per
temperature\tabularnewline
\bottomrule
\end{longtable}

The full example takes a few minutes on a laptop. For a faster check,
reduce \texttt{Trange1} to two temperatures.

\subsection{Programmatic use}\label{programmatic-use}

\begin{lstlisting}
from mag_anal.Anal_Tau import CAL_RELAX

tau = CAL_RELAX("inpfile/CFPs_v1.yml")   # run from examples/relaxation
tau.kernel()

# Orbach rate-determining analysis for the first temperature
tau.TP[0].Mod_Analyse_OB()
\end{lstlisting}

\texttt{tau.TP} is a list of \texttt{TPs\_Analyze} objects, one per
temperature.

\begin{center}\rule{0.5\linewidth}{0.4pt}\end{center}

\section{Full SPC workflow}\label{full-spc-workflow}

\subsection{Step 1 --- geometry optimisation and
phonons}\label{step-1-geometry-optimisation-and-phonons}

Prepare in \texttt{SPC\_Tool/inpfile/}:

\begin{itemize}
\tightlist
\item
  \texttt{POSCAR} (or a \texttt{.cif} for \texttt{gen\_POSCAR.py});
\item
  \texttt{opt\_INCAR}, \texttt{freq\_INCAR};
\item
  \texttt{job\_vasp.sh} adapted to your cluster.
\end{itemize}

Run, in \texttt{SPC\_Tool/opt\_freq/vaspsh/}:

\begin{lstlisting}
bash POSCAR.sh     # optional: build POSCAR from a structure
bash opt.sh        # VASP geometry optimisation
bash freq.sh       # VASP frequencies; calls Select_group.py
\end{lstlisting}

\texttt{freq.sh} checks that the optimisation reached
\texttt{reached\ required\ accuracy}, then creates \texttt{VASP/freq/},
runs \texttt{Select\_group.py} and copies \texttt{freq\_POSCAR},
\texttt{eq\_opt.xyz}, \texttt{dpxs} and \texttt{Atom\_list} into
\texttt{inpfile/}.

\texttt{Select\_group.py} options:

\begin{longtable}[]{@{}>{\raggedright\arraybackslash}p{0.30\linewidth} >{\raggedright\arraybackslash}p{0.66\linewidth}@{}}
\toprule
Option & Meaning\tabularnewline
\midrule
\endhead
\texttt{-\/-cry} & crystal structure file\tabularnewline
\texttt{-\/-ct} & centre element (e.g. \texttt{Dy})\tabularnewline
\texttt{-\/-num} & centre atom index (optional)\tabularnewline
\texttt{-\/-out} & output folder\tabularnewline
\texttt{-f}, \texttt{-\/-freq} & also write a
\texttt{Selective\ dynamics} \texttt{freq\_POSCAR}\tabularnewline
\bottomrule
\end{longtable}

\subsection{Step 2 --- displaced
structures}\label{step-2-displaced-structures}

In \texttt{SPC\_Tool/Per\_B\_kp/Pyscfsh/}:

\begin{lstlisting}
bash run1_stru.sh
bash run2_Regr.sh   # only for the periodic-cell path
\end{lstlisting}

\texttt{run1\_stru.sh} calls either:

\begin{lstlisting}
python3 gen_structures_xyz.py --pos ../../inpfile/eq_opt.xyz -np 10 -mdp 0.02
# or
python3 gen_structures_cell.py --out ../../inpfile/freq.log --pos ../../inpfile/opt_CONTCAR -np 10 -mdp 0.2
\end{lstlisting}

\begin{longtable}[]{@{}>{\raggedright\arraybackslash}p{0.30\linewidth} >{\raggedright\arraybackslash}p{0.66\linewidth}@{}}
\toprule
Script & Options\tabularnewline
\midrule
\endhead
\texttt{gen\_structures\_xyz.py} & \texttt{-\/-pos},
\texttt{-np/-\/-npts}, \texttt{-mdp/-\/-maxdp}\tabularnewline
\texttt{gen\_structures\_cell.py} & \texttt{-\/-out} (VASP frequency
file), \texttt{-\/-pos}, \texttt{-np}, \texttt{-mdp},
\texttt{-v}\tabularnewline
\texttt{Replace\_group.py} & \texttt{-i/-\/-inp}, \texttt{-o/-\/-out},
\texttt{-\/-inpfile}\tabularnewline
\bottomrule
\end{longtable}

\texttt{gen\_structures\_*} writes
\texttt{perb/\textless{}mode\textgreater{}/\textless{}point\textgreater{}/mol\_\textless{}mode\textgreater{}\_\textless{}point\textgreater{}.xyz}.
\texttt{Replace\_group.py} extracts the molecular group from the
periodic cell using \texttt{inpfile/Atom\_list} and
\texttt{inpfile/dpxs}.

\subsection{Step 3 --- electronic-structure
inputs}\label{step-3-electronic-structure-inputs}

\begin{lstlisting}
bash run3_inp.sh
\end{lstlisting}

\texttt{run3\_inp.sh} fills the template for the equilibrium structure
and every displacement:

\begin{lstlisting}
python insert_xyz.py --tmp ../../inpfile/dmet_tmp.py.tmpl \
    --xyz mol_<mode>_<point>.xyz -o dmet_cassiso.py -c <charge> -t <title>
\end{lstlisting}

\texttt{insert\_xyz.py} options:

\begin{longtable}[]{@{}>{\raggedright\arraybackslash}p{0.30\linewidth} >{\raggedright\arraybackslash}p{0.66\linewidth}@{}}
\toprule
Option & Meaning\tabularnewline
\midrule
\endhead
\texttt{-\/-tmp} & template file (\texttt{.tmpl})\tabularnewline
\texttt{-\/-xyz} & structure file\tabularnewline
\texttt{-o/-\/-out} & output Python input\tabularnewline
\texttt{-c/-\/-charge} & total charge\tabularnewline
\texttt{-t/-\/-title} & job title\tabularnewline
\bottomrule
\end{longtable}

The resulting \texttt{dmet\_cassiso.py} still needs a \textbf{complete
basis block}. The template contains Dy/C/H entries as an example; add
every element present in your system.

For OpenMolcas, use \texttt{SPC\_Tool/inpfile/molcas\_cas.in} and the
Molcas job templates instead.

\subsection{Step 4 --- run the electronic-structure
jobs}\label{step-4-run-the-electronic-structure-jobs}

The generated layout is:

\begin{lstlisting}
perb/
|-- eq/
|   |-- dmet_cassiso.py
|   |-- jobeq.sh
|-- <mode>/<point>/
    |-- dmet_cassiso.py
    |-- pyjob.sh
\end{lstlisting}

Submit \texttt{jobeq.sh} first, then the displaced jobs.
\texttt{SPC\_Tool/sh/} contains templates:

\begin{longtable}[]{@{}>{\raggedright\arraybackslash}p{0.30\linewidth} >{\raggedright\arraybackslash}p{0.66\linewidth}@{}}
\toprule
Script & Purpose\tabularnewline
\midrule
\endhead
\texttt{job\_pyscf.sh} & one displaced PySCF job\tabularnewline
\texttt{jobeq\_pyscf.sh} & equilibrium PySCF job\tabularnewline
\texttt{job\_molcas.sh}, \texttt{jobeq\_molcas.sh} & OpenMolcas
equivalents\tabularnewline
\texttt{check\_pyscf.sh}, \texttt{check\_molcas.sh} & check completion
and trim output\tabularnewline
\texttt{clean\_pyscf.sh} & remove checkpoint files\tabularnewline
\texttt{yhjob\_pyscf.sh}, \texttt{yhjob\_molcas.sh} & batch
submission\tabularnewline
\bottomrule
\end{longtable}

These files are \textbf{templates}; edit partition/queue, e-mail,
modules and \texttt{VASP\_BIN}. See \S10.

\subsection{Step 5 --- extract CFPs}\label{step-5-extract-cfps}

\begin{lstlisting}
bash run4_CFPs.sh
\end{lstlisting}

which is equivalent to:

\begin{lstlisting}
python3 Extract_CFPs.py -el perb/eq/pyout -nd 201 -np 10
\end{lstlisting}

\begin{longtable}[]{@{}>{\raggedright\arraybackslash}p{0.30\linewidth} >{\raggedright\arraybackslash}p{0.66\linewidth}@{}}
\toprule
Option & Meaning\tabularnewline
\midrule
\endhead
\texttt{-el/-\/-eqlog} & equilibrium magnetic output\tabularnewline
\texttt{-nd/-\/-nmds} & number of vibrational modes\tabularnewline
\texttt{-np/-\/-npts} & number of displacement points\tabularnewline
\texttt{-\/-prefix} & directory containing \texttt{perb} (default
\texttt{perb})\tabularnewline
\texttt{-\/-stdout} & per-point output name (default
\texttt{stdout1})\tabularnewline
\texttt{-\/-out} & output CFP file (default
\texttt{CFPs.out})\tabularnewline
\texttt{-v} & verbose\tabularnewline
\bottomrule
\end{longtable}

The result \texttt{CFPs.out} is the \texttt{CFPs\_filename} consumed by
\texttt{mag\_anal}.

\subsection{Step 6 --- relaxation
time}\label{step-6-relaxation-time}

Create a YAML file pointing to \texttt{CFPs.out}, the equilibrium
output, \texttt{freq.log}, \texttt{POSCAR}, \texttt{Atom\_list} and the
tensor-operator folder, then run:

\begin{lstlisting}
lamp-tau --yml my_system.yml
\end{lstlisting}

See \S6 for all keys and \S8 for outputs.

\begin{center}\rule{0.5\linewidth}{0.4pt}\end{center}

\section{YAML reference}\label{yaml-reference}

A complete example is \texttt{examples/relaxation/CFPs\_v1.yml}.

\subsection{Required keys}\label{required-keys}

\begin{longtable}[]{@{}>{\raggedright\arraybackslash}p{0.26\linewidth} >{\raggedright\arraybackslash}p{0.14\linewidth} >{\raggedright\arraybackslash}p{0.58\linewidth}@{}}
\toprule
Key\strut
 & Type\strut
 & Meaning\strut
\tabularnewline
\midrule
\endhead
\texttt{CFPs\_filename}\strut
 & path\strut
 & \texttt{CFPs.out} from \texttt{Extract\_CFPs.py}\strut
\tabularnewline
\texttt{mag\_filename}\strut
 & path\strut
 & equilibrium electronic-structure output\strut
\tabularnewline
\texttt{Ab\_Type}\strut
 & int\strut
 & \texttt{0} Molcas, \texttt{1} OpenMolcas, \texttt{3} PySCF\strut
\tabularnewline
\texttt{freq\_filename}\strut
 & path\strut
 & VASP frequency output\strut
\tabularnewline
\texttt{POSCAR}\strut
 & path\strut
 & equilibrium structure\strut
\tabularnewline
\texttt{TO\_foldname}\strut
 & path\strut
 & tensor-operator folder\strut
\tabularnewline
\texttt{VP\_filename}\strut
 & path\strut
 & spin-phonon matrix output/prefix\strut
\tabularnewline
\texttt{TPs\_foldname}\strut
 & path\strut
 & cache folder for transition matrices\strut
\tabularnewline
\texttt{Out\_foldname}\strut
 & path\strut
 & folder for \(\tau\)(T) outputs\strut
\tabularnewline
\texttt{CFPs\_X}\strut
 & list\strut
 & displacement grid; must contain \texttt{0}\strut
\tabularnewline
\texttt{Correction\_slope}, \texttt{Corrected\_intercept}\strut
 & float\strut
 & frequency correction \texttt{\(\nu\)\ →\ a·\(\nu\)\ +\ b}\strut
\tabularnewline
\texttt{Fitting\_order}\strut
 & int\strut
 & polynomial order of the CFP fit\strut
\tabularnewline
\texttt{Linewidth\_parameter}\strut
 & float\strut
 & broadening (cm⁻¹)\strut
\tabularnewline
\texttt{If\_Data}\strut
 & bool\strut
 & overlay experimental \(\tau\)(T)\strut
\tabularnewline
\texttt{If\_Perb}\strut
 & bool\strut
 & apply the perturbation treatment\strut
\tabularnewline
\texttt{Perb\_T}\strut
 & float\strut
 & perturbation field (T)\strut
\tabularnewline
\texttt{g\_z}\strut
 & float\strut
 & ground-state \texttt{g\_z}\strut
\tabularnewline
\texttt{If\_uniformization}\strut
 & bool\strut
 & use the uniformization method\strut
\tabularnewline
\texttt{CFPs\_Type}\strut
 & int\strut
 & see \S6.2\strut
\tabularnewline
\texttt{If\_get\_picture\_Tt}\strut
 & bool\strut
 & write \(\tau\)(T) figures\strut
\tabularnewline
\texttt{If\_Orbach}, \texttt{If\_Raman}\strut
 & bool\strut
 & enable relaxation channels\strut
\tabularnewline
\bottomrule
\end{longtable}

\subsection{\texorpdfstring{\texttt{CFPs\_Type}}{6.2 CFPs\_Type}}\label{cfps_type}

\begin{longtable}[]{@{}>{\raggedright\arraybackslash}p{0.26\linewidth} >{\raggedright\arraybackslash}p{0.14\linewidth} >{\raggedright\arraybackslash}p{0.58\linewidth}@{}}
\toprule
Value & Meaning & Extra keys\tabularnewline
\midrule
\endhead
\texttt{0} & normal-mode CFPs & none\tabularnewline
\texttt{1} & Cartesian (xyz) CFPs & \texttt{Atom\_list}\tabularnewline
\texttt{2} & xyz + spatial reduction & \texttt{Molfile}, plus
\texttt{Atom\_radius} \textbf{or} \texttt{Atom\_label}\tabularnewline
\texttt{3} & Lorentz random CFPs & \texttt{llzA},
\texttt{llzw}\tabularnewline
\bottomrule
\end{longtable}

For \texttt{CFPs\_Type=1} and \texttt{2}, \texttt{Atom\_list} points to
the atom indices retained by the spatial-truncation selection. The
default is \texttt{inpfile/Atom\_list}; specify it explicitly for
portable examples.

\subsection{Optional keys}\label{optional-keys}

\begin{longtable}[]{@{}>{\raggedright\arraybackslash}p{0.30\linewidth} >{\raggedright\arraybackslash}p{0.66\linewidth}@{}}
\toprule
Key\strut
 & Meaning\strut
\tabularnewline
\midrule
\endhead
\texttt{Atom\_list}\strut
 & atom indices for the mode projection (default
\texttt{inpfile/Atom\_list})\strut
\tabularnewline
\texttt{First\_list}\strut
 & first-shell indices for \texttt{CFPs\_Type=2}\strut
\tabularnewline
\texttt{Molfile}, \texttt{ct\_idx}, \texttt{Atom\_radius},
\texttt{Atom\_label}\strut
 & spatial-reduction options\strut
\tabularnewline
\texttt{llzA}, \texttt{llzw}\strut
 & Lorentz random-CFP parameters\strut
\tabularnewline
\texttt{If\_select\_sigma}, \texttt{select\_sigma}\strut
 & restrict modes to a frequency window (cm⁻¹)\strut
\tabularnewline
\texttt{If\_T\_step}, \texttt{nTrange}, \texttt{Trange1..n}\strut
 & temperature grid\strut
\tabularnewline
\texttt{ndata}, \texttt{Tr1..}, \texttt{logtime1..},
\texttt{label1..}\strut
 & experimental data overlay\strut
\tabularnewline
\bottomrule
\end{longtable}

\subsection{Temperature grid}\label{temperature-grid}

\begin{itemize}
\tightlist
\item
  \texttt{If\_T\_step:\ True}: each \texttt{TrangeN} is
  \texttt{min,max,step}.
\item
  \texttt{If\_T\_step:\ False}: each \texttt{TrangeN} is
  \texttt{min,max,count}.
\end{itemize}

Example:

\begin{lstlisting}
If_T_step: True
nTrange: 1
Trange1: 10,50,5      # 10, 15, ..., 50 K
\end{lstlisting}

\subsection{Experimental overlay}\label{experimental-overlay}

With \texttt{If\_Data:\ True}, provide:

\begin{lstlisting}
ndata: 1
Tr1: [ ... ]          # 1/T
logtime1: [ ... ]     # log10(tau/s)
label1: Exp
\end{lstlisting}

\begin{center}\rule{0.5\linewidth}{0.4pt}\end{center}

\section{File formats}\label{file-formats}

\subsection{\texorpdfstring{\texttt{CFPs.out}}{7.1 CFPs.out}}\label{cfps.out}

\begin{lstlisting}
The Perturbation CFPs
<blank line>
Total modes 201
Total points 10
Mode 1 Point 1
<5 numbers>
<9 numbers>
<13 numbers>
...
\end{lstlisting}

The blank line is optional in recent versions (the reader searches for
\texttt{Total\ modes}/\texttt{Total\ points}), but
\texttt{Extract\_CFPs.py} writes it.

\subsection{\texorpdfstring{\texttt{freq.log}}{7.2 freq.log}}\label{freq.log}

VASP frequency output with blocks:

\begin{lstlisting}
   1 f  =  102.701176 THz ... 3425.742380 cm-1 ...
             X         Y         Z           dx          dy          dz
      ...
\end{lstlisting}

The mode order must match \texttt{CFPs.out}.

\subsection{\texorpdfstring{\texttt{eqpyout}}{7.3 eqpyout}}\label{eqpyout}

PySCF output containing the
\texttt{Rotated\ Crystal\ Field\ Parameters:} table and spin-orbit
eigenstates. For Molcas/OpenMolcas use \texttt{Ab\_Type=0/1}.

\subsection{\texorpdfstring{\texttt{.eso}}{7.4 .eso}}\label{eso}

\texttt{TO\_foldname} contains
\texttt{\textless{}Nstate\textgreater{}\_real.eso} and
\texttt{\textless{}Nstate\textgreater{}\_imag.eso}:

\begin{lstlisting}
k = 2,q = -2
<16 x 16 real numbers>
...
\end{lstlisting}

The shipped \texttt{examples/relaxation/ESO\_TTO/} files are for 16
states (Dy³⁺).

\subsection{Templates}\label{templates}

\texttt{SPC\_Tool/inpfile/*.tmpl} are \textbf{not runnable}. They
contain three placeholder lines:

\begin{lstlisting}
Molecular_Name
Molecular_Charge
Calculate_Title
\end{lstlisting}

\texttt{insert\_xyz.py} replaces these lines. Complete the basis block
before running.

\begin{center}\rule{0.5\linewidth}{0.4pt}\end{center}

\section{Output files}\label{output-files}

\subsection{\texorpdfstring{\texttt{Out\_foldname}}{8.1 Out\_foldname}}\label{out_foldname}

\begin{longtable}[]{@{}>{\raggedright\arraybackslash}p{0.30\linewidth} >{\raggedright\arraybackslash}p{0.66\linewidth}@{}}
\toprule
File\strut
 & Content\strut
\tabularnewline
\midrule
\endhead
\texttt{OB\_v1.dat}\strut
 & columns \texttt{1/T}, \texttt{-log10(\(\tau\)/s)}, standard deviation\strut
\tabularnewline
\texttt{RM\_v1.dat}\strut
 & same for the Raman channel\strut
\tabularnewline
\texttt{TPs1.png}\strut
 & \texttt{log10(\(\tau\)/s)} vs \texttt{1/T}\strut
\tabularnewline
\texttt{TPs2.png}\strut
 & comparison figure\strut
\tabularnewline
\texttt{uniform.dat}\strut
 & uniformization results (\texttt{1/T}, \texttt{-log10(rate)} for three
estimates)\strut
\tabularnewline
\texttt{OBTrans.eps}, \texttt{RMTrans.eps}, \texttt{RMPie.eps}\strut
 & mode-resolved analysis (after
\texttt{Mod\_Analyse\_OB}/\texttt{RM})\strut
\tabularnewline
\bottomrule
\end{longtable}

The Y-axis of \texttt{TPs1.png} is \texttt{log10(\(\tau\)/s)}; column 2 of
\texttt{OB\_v1.dat} is its negative (\texttt{-log10(\(\tau\)/s)}). Negate
column 2 to recover the plotted quantity.

\subsection{\texorpdfstring{\texttt{TPs\_foldname}}{8.2 TPs\_foldname}}\label{tps_foldname}

\begin{longtable}[]{@{}>{\raggedright\arraybackslash}p{0.30\linewidth} >{\raggedright\arraybackslash}p{0.66\linewidth}@{}}
\toprule
File & Content\tabularnewline
\midrule
\endhead
\texttt{OB\textless{}T\textgreater{}.tps} & Orbach transition matrix at
\texttt{T} K (text)\tabularnewline
\texttt{OB\textless{}T\textgreater{}.tp} & summed Orbach
matrix\tabularnewline
\texttt{LM\textless{}T\textgreater{}.tps} & Raman transition
matrix\tabularnewline
\texttt{LM\textless{}T\textgreater{}.tp} & summed Raman
matrix\tabularnewline
\bottomrule
\end{longtable}

If a cached file exists, it is read instead of recomputed. Delete it to
force recalculation. The
\texttt{If\_select\_sigma}/\texttt{select\_sigma} header is checked by
\texttt{check\_select\_sigma}.

\subsection{\texorpdfstring{\texttt{VP\_filename}}{8.3 VP\_filename}}\label{vp_filename}

The spin-phonon matrix. A \texttt{.dat} file is written first and
renamed on completion; an interrupted run leaves \texttt{*.dat} and can
be removed safely.

\subsection{\texorpdfstring{\texttt{psh.py}}{8.4 psh.py}}\label{psh.py}

\texttt{PSH} writes \texttt{SOC.txt}, \texttt{Lx/Ly/Lz.txt},
\texttt{Sx/Sy/Sz.txt}, \texttt{eigval.txt} and \texttt{eig\_*.txt} to
\texttt{outdir} (default: current directory).

\begin{center}\rule{0.5\linewidth}{0.4pt}\end{center}

\section{Python API}\label{python-api}

\subsection{\texorpdfstring{\texttt{CAL\_RELAX}}{9.1 CAL\_RELAX}}\label{cal_relax}

\begin{lstlisting}
from mag_anal.Anal_Tau import CAL_RELAX

tau = CAL_RELAX("my_system.yml")
tau.kernel()
\end{lstlisting}

\texttt{kernel()} calls \texttt{Read\_input}, \texttt{Read\_freq},
\texttt{Read\_CFPs}, \texttt{Read\_Static\_effH},
\texttt{Analyse\_Relax}.

Useful attributes after \texttt{kernel()}:

\begin{longtable}[]{@{}>{\raggedright\arraybackslash}p{0.30\linewidth} >{\raggedright\arraybackslash}p{0.66\linewidth}@{}}
\toprule
Attribute & Meaning\tabularnewline
\midrule
\endhead
\texttt{tau.T\_list} & temperature grid\tabularnewline
\texttt{tau.TP} & list of \texttt{TPs\_Analyze} (one per
T)\tabularnewline
\texttt{tau.CFPs} & \texttt{CFPs\_Analyze} object\tabularnewline
\texttt{tau.Eigen} & \texttt{Eigenstates} object\tabularnewline
\texttt{tau.Out\_foldname} & output folder\tabularnewline
\bottomrule
\end{longtable}

\subsection{\texorpdfstring{\texttt{CFPs\_Analyze}}{9.2 CFPs\_Analyze}}\label{cfps_analyze}

\begin{lstlisting}
from mag_anal.Anal_Tau import CFPs_Analyze

cfps = CFPs_Analyze("CFPs.out", CFPs_X, atom_list="Atom_list")
cfps.propare_CFPs(mol, CFPs_Type=1, order=1)
\end{lstlisting}

Methods: \texttt{read\_CFPs}, \texttt{read\_CFPs\_xyz},
\texttt{read\_gtensors}, \texttt{propare\_CFPs}, \texttt{gen\_slope},
\texttt{gen\_slope\_xyz}, \texttt{gen\_slope\_reduced},
\texttt{transB\_sl}, \texttt{transB\_sl\_select}, \texttt{write\_CFPs}.

\subsection{\texorpdfstring{\texttt{Eigenstates}}{9.3 Eigenstates}}\label{eigenstates}

Reads energies/wavefunctions from PySCF (\texttt{read\_energies\_Pyscf},
\texttt{read\_wavefunction\_Pyscf}), Molcas (\texttt{\_M}) or OpenMolcas
(\texttt{\_OM}), builds the tensor-operator matrices and the spin-phonon
matrix.

\subsection{\texorpdfstring{\texttt{TPs\_Analyze}}{9.4 TPs\_Analyze}}\label{tps_analyze}

\begin{longtable}[]{@{}>{\raggedright\arraybackslash}p{0.30\linewidth} >{\raggedright\arraybackslash}p{0.66\linewidth}@{}}
\toprule
Method & Purpose\tabularnewline
\midrule
\endhead
\texttt{read\_OBTPs}, \texttt{read\_RMTPs} & read cached
matrices\tabularnewline
\texttt{sum\_OBRMTPs} & sum channels\tabularnewline
\texttt{Mod\_Analyse\_OB}, \texttt{Mod\_Analyse\_RM} &
rate-determining-step analysis\tabularnewline
\texttt{save\_TP}, \texttt{save\_TPP} & write matrices\tabularnewline
\bottomrule
\end{longtable}

\subsection{\texorpdfstring{\texttt{PSH} (pseudospin
Hamiltonian)}{9.5 PSH (pseudospin Hamiltonian)}}\label{psh-pseudospin-hamiltonian}

\begin{lstlisting}
from mag_anal.psh import PSH

psh = PSH(mysiso, Nstate=16, lmax=12, TFO=1, TFB=0, ROT=True, outdir="psh_out")
psh.kernel()
\end{lstlisting}

\begin{longtable}[]{@{}>{\raggedright\arraybackslash}p{0.30\linewidth} >{\raggedright\arraybackslash}p{0.66\linewidth}@{}}
\toprule
Argument & Meaning\tabularnewline
\midrule
\endhead
\texttt{SISO} & \texttt{embed\_sim.siso.SISO} object\tabularnewline
\texttt{Nstate} & number of \texttt{\textbar{}J\ M\textgreater{}}
states\tabularnewline
\texttt{lmax} & maximum tensor rank\tabularnewline
\texttt{TFO} & \texttt{0} EasySpin ESO, \texttt{1} tesseral tensor
operators (Molforge style)\tabularnewline
\texttt{TFB} & \texttt{0} least squares, \texttt{1}
projection\tabularnewline
\texttt{ROT} & rotate CFPs to the easy axis\tabularnewline
\texttt{Angle} & explicit ZYZ Euler angles; \texttt{None} uses the easy
axis\tabularnewline
\texttt{outdir} & directory for diagnostic text files\tabularnewline
\bottomrule
\end{longtable}

\texttt{psh.kernel()} runs \texttt{gen\_ESO}, \texttt{get\_E\_J},
\texttt{get\_g\_tensor} and \texttt{get\_Bkq}/\texttt{get\_rot\_Bkq}.

\subsection{\texorpdfstring{\texttt{gen\_structures} /
\texttt{Select\_group}}{9.6 gen\_structures / Select\_group}}\label{gen_structures-select_group}

\begin{lstlisting}
from mag_anal.gen_structures import mol_freq
from mag_anal.Select_group import choose_group
\end{lstlisting}

\texttt{mol\_freq(outfile,\ posfile)} reads VASP phonons;
\texttt{choose\_group(cryfile,\ center,\ outfold,\ r1,\ r2)} selects a
molecular group.

\begin{center}\rule{0.5\linewidth}{0.4pt}\end{center}

\section{HPC notes}\label{hpc-notes}

The files under \texttt{SPC\_Tool/sh/} are templates. Before submitting:

\begin{itemize}
\tightlist
\item
  replace \texttt{YOUR\_PARTITION} with your partition/queue;
\item
  replace \texttt{YOUR\_EMAIL} or remove the mail line;
\item
  set \texttt{VASP\_BIN} instead of a hard-coded path;
\item
  replace \texttt{module\ load\ ...} /
  \texttt{\#\ source\ activate\ ...} with the commands that provide
  PySCF, LAMP-emb and ASE;
\item
  make \texttt{num\_freq}, \texttt{num\_pt} and \texttt{name} consistent
  across \texttt{run1\_stru.sh}, \texttt{run2\_Regr.sh},
  \texttt{run3\_inp.sh}, \texttt{run4\_CFPs.sh} and the job scripts.
\end{itemize}

Slurm example:

\begin{lstlisting}
#!/bin/bash
#SBATCH -o stdout
#SBATCH --partition=YOUR_PARTITION
#SBATCH -J lamp_tau
#SBATCH --nodes=1
#SBATCH --ntasks-per-node=28
#SBATCH --time=100:00:00

module purge
module load pyscf
export OMP_NUM_THREADS=7
export PYTHONPATH=/path/to/LAMP_emb:$PYTHONPATH
python dmet_cassiso.py
\end{lstlisting}

PBS example:

\begin{lstlisting}
#!/bin/bash
#PBS -N lamp_tau
#PBS -q YOUR_QUEUE
#PBS -l nodes=1:ppn=28
#PBS -l walltime=100:00:00
#PBS -o stdout
#PBS -j oe

cd "$PBS_O_WORKDIR"
python dmet_cassiso.py
\end{lstlisting}

The original scripts also used Tianhe \texttt{yhbatch}; replace it with
the current scheduler commands.

\begin{center}\rule{0.5\linewidth}{0.4pt}\end{center}

\section{Troubleshooting}\label{troubleshooting}

\begin{longtable}[]{@{}>{\raggedright\arraybackslash}p{0.30\linewidth} >{\raggedright\arraybackslash}p{0.66\linewidth}@{}}
\toprule
Symptom\strut
 & Cause / fix\strut
\tabularnewline
\midrule
\endhead
\texttt{ModuleNotFoundError:\ embed\_sim}\strut
 & Install LAMP-emb and add it to \texttt{PYTHONPATH}.\strut
\tabularnewline
\texttt{ModuleNotFoundError:\ mag\_anal}\strut
 & Install with \texttt{pip\ install\ -e\ .} or run from the repository
root.\strut
\tabularnewline
\texttt{FileNotFoundError:\ .../Atom\_list}\strut
 & Set \texttt{Atom\_list} in the YAML or create
\texttt{inpfile/Atom\_list}.\strut
\tabularnewline
\texttt{NotImplementedError:\ Ab\_Type=2}\strut
 & ORCA output is not implemented; use \texttt{Ab\_Type=0/1/3}.\strut
\tabularnewline
\texttt{FileNotFoundError:\ CFPs.out}\strut
 & Run \texttt{run4\_CFPs.sh} first.\strut
\tabularnewline
\texttt{KeyError:\ Missing\ required\ key}\strut
 & Add the missing key to the YAML (the message lists it).\strut
\tabularnewline
\texttt{ValueError:\ Could\ not\ find\ \textquotesingle{}Total\ modes\textquotesingle{}}\strut
 & The CFP file is not in the expected format; see \S7.1.\strut
\tabularnewline
Plots do not open on a server\strut
 & Set \texttt{MPLBACKEND=Agg}.\strut
\tabularnewline
\texttt{TPs} values look stale\strut
 & Delete the cached
\texttt{OB\textless{}T\textgreater{}.tps}/\texttt{LM\textless{}T\textgreater{}.tps}
and rerun.\strut
\tabularnewline
Shell script does nothing\strut
 & The scheduler header/modules still belong to the author's cluster; edit
them (\S10).\strut
\tabularnewline
\texttt{vasp\_std:\ command\ not\ found}\strut
 & Set \texttt{export\ VASP\_BIN=/path/to/vasp\_std}.\strut
\tabularnewline
\bottomrule
\end{longtable}

\begin{center}\rule{0.5\linewidth}{0.4pt}\end{center}

\section{Validation}\label{validation}

The shipped DyAC example was used as a smoke test:

\begin{itemize}
\tightlist
\item
  \texttt{python\ -m\ compileall\ mag\_anal\ SPC\_Tool} passes;
\item
  from \texttt{examples/relaxation},
  \texttt{python\ -m\ mag\_anal.Anal\_Tau\ -\/-yml\ inpfile/CFPs\_v1.yml}
  completes;
\item
  the 201-mode, 10-50 K run reports \texttt{Ueff\ ≈\ 190\ K} and
  \texttt{log(\(\tau\)/s)\ ≈\ -11.3};
\item
  \texttt{tests/test\_smoke.py} checks that \texttt{OB\_v1.dat} is
  produced (uses two temperatures).
\end{itemize}

Run:

\begin{lstlisting}
MPLBACKEND=Agg python -m pytest tests -q
\end{lstlisting}

\begin{center}\rule{0.5\linewidth}{0.4pt}\end{center}

\section{Citation, license and
support}\label{citation-license-and-support}

If you use LAMP-tau, please cite the manuscript and the archived
release:

\begin{itemize}
\tightlist
\item
  manuscript: \emph{Accelerating ab initio spin-phonon relaxation
  simulation of single-ion magnets by quantum embedding and spatial
  truncation}, J. Chem. Phys. (2026);
\item
  software: see \texttt{CITATION.cff}; the Zenodo DOI is added after the
  first release.
\end{itemize}

License: MIT (see \texttt{LICENSE}). Third-party programs keep their own
licenses:

\begin{itemize}
\tightlist
\item
  PySCF --- Apache-2.0;
\item
  LAMP-emb --- not redistributed here; install it separately;
\item
  VASP, OpenMolcas, ORCA --- their own licenses;
\item
  ASE, NumPy, SciPy, Matplotlib, SymPy, prettytable, PyYAML --- their
  own licenses.
\end{itemize}

For questions, open an issue on GitHub. For scientific questions,
contact the corresponding authors of the paper.

\end{document}
